\section{The Data\label{sec:thedata}}
\subsection{BKK's Two Papers}
Backus, Kehoe and Kydland examine international business cycles in two papers: ``International Real Business Cycles'' (1992) and ``International Business Cycles: Theory and Evidence'' (1993). Both report output and consumption correlations between the United States and other countries, but their data sources, country samples and observation periods differ. These choices lead to meaningful differences in their results.

BKK (1993) use the same sample period for every country and variable: 1970:1 to 1990:2. Their data come from the OECD's \textit{Quarterly National Accounts} and \textit{Main Economic Indicators}.\footnote{BKK (1993) also use the IMF's \textit{International Financial Statistics} for exchange-rate data, which are outside the scope of this paper.} They examine 11 countries, two fewer than BKK (1992); both samples include a European aggregate. BKK (1992) use the IMF's \textit{International Financial Statistics}, with sample periods that vary by country. The longest series covers the United States from 1960:1 to 1990:2, while the shortest covers the European aggregate from 1970:1 to 1986:4.

The differences in sources and sample periods matter. The Australia--US output correlation differs by 0.26 between the papers. The France--US consumption correlation differs by 0.40, rising from -0.18 in BKK (1992) to 0.22 in BKK (1993). These discrepancies motivate our approach to data collection: because we are reassessing BKK (1993), we use sources and measures as close as possible to theirs to make the comparison meaningful.


Despite these differences, the papers share several central findings. Output correlations are positive for all countries except South Africa in BKK (1992), and consumption is much less highly correlated across countries than the model predicts. Some consumption correlations with the United States are even negative: Australia in BKK (1993), and Finland, France and South Africa in BKK (1992).

Throughout the remainder of this paper, BKK refers to ``International Business Cycles: Theory and Evidence'' (1993).

\subsection{Our Data}

Our sample contains the same countries as BKK's: Australia (AUS), Austria (AUT), Canada (CAN), France (FRA), Germany (DEU), Italy (ITA), Japan (JPN), Switzerland (CHE), the United Kingdom (GBR), the United States (USA), and a European aggregate (EU15).\footnote{Our European aggregate comprises fifteen countries: Austria, Belgium, Denmark, Finland, France, Germany, Greece, Ireland, Italy, Luxembourg, the Netherlands, Portugal, Spain, Sweden and the United Kingdom. BKK do not specify the composition of their European aggregate.}

We also draw on the OECD's \textit{Quarterly National Accounts} and \textit{Main Economic Indicators}. BKK provide limited information about their measurement choices, stating only that output, consumption, fixed investment and government purchases are measured in real terms, while the net exports-to-output ratio is calculated at current prices. Of the 20 potentially relevant OECD measures, we select VPVOBARSA\footnote{VPVOBARSA denotes seasonally adjusted volume estimates at constant 2010 prices, converted using 2010 purchasing power parities and expressed in US dollars at annual levels.} for output, consumption, investment and government spending. We use CPCARSA,\footnote{CPCARSA denotes seasonally adjusted values at current prices and current purchasing power parities, expressed in US dollars at annual levels.} its current-price counterpart, to calculate the net exports-to-output ratio. Appendix~\ref{measure} explains these choices.

Following BKK, we calculate Solow residuals from $\log z = \log y - [\theta \log k + (1-\theta)\log n]$, omitting the capital term $\theta\log k$.

We construct the civilian employment series by combining the OECD's \textit{Short-Term Labour Market Statistics} and \textit{Main Economic Indicators}, as described in Appendix~\ref{app:employment}. \Cref{tab:variables} summarizes the variables and their OECD codes and measures.

\begin{table}[htbp]
  \centering

\resizebox{\textwidth}{!}{\begin{tabular}{llcc}
\toprule
          & \multicolumn{1}{c}{Name}  & Code  & Measure \\
\cmidrule(lr){2-4}
    y     & Gross domestic product - expenditure approach & B1\_GE & VPVOBARSA \\
    c     & Private final consumption expenditure & P31S14\_S15 & VPVOBARSA \\
    x     & Gross fixed capital formation & P51   & VPVOBARSA \\
    g     &  General government final consumption expenditure & P3S13 & VPVOBARSA \\
 n     &  Employed population, Aged 15 and over, All persons & LFEMTTTT & STSA \\
\cmidrule(lr){2-4}
  \multirow{3}{*}{nx}   & \multicolumn{1}{l}{Exports of goods and services} & \multicolumn{1}{c}{P6}    & \multicolumn{1}{c}{CPCARSA} \\
               & \multicolumn{1}{l}{Imports of goods and services} & \multicolumn{1}{c}{P7}    & \multicolumn{1}{c}{CPCARSA} \\
              & \multicolumn{1}{l}{Gross domestic product - expenditure approach} & \multicolumn{1}{c}{B1\_GE} & \multicolumn{1}{c}{CPCARSA} \\\cmidrule(lr){2-4}
\bottomrule
    \end{tabular}}%
  \caption{OECD Variable Definitions and Measures}
  \label{tab:variables}
\end{table}%


We extend BKK's sample both forward, adding 23 years of more recent data, and backward, including earlier observations available in the OECD database. \Cref{tab:timespan} reports the coverage of each variable and country. Most series run from 1960:1 to 2015:1, although some US and UK series begin in 1947:1 and 1955:1, respectively. \cite{Backus:1992aaa} provide estimates for longer periods, including years before the Second World War. These earlier estimates are less reliable than the post-1945 data and are not considered here.

\generatedtable{table_3_timespan.tex}{Time Series Coverage by Country}{tab:timespan}

\subsection{Hodrick-Prescott Filter}

Analyzing international comovements requires a definition of the business cycle. Economic time series can be decomposed into a slowly varying growth trend and a more rapidly changing cyclical component. Neither component is directly observable; only the resulting values of output, investment and other variables are observed. The researcher must therefore choose a method for estimating the trend and identifying deviations from it.

For comparability with BKK, we use the Hodrick-Prescott (HP) filter described in \cite{Hodrick:1981aa}. The filter emphasizes medium- and high-frequency movements and ``can be thought of as removing a smooth trend line (like one might draw freehand) from the data'' (\cite{Backus:1992aaa}). According to \cite{marianne1}, the resulting cyclical component is stationary if the input series is integrated of order four or less. \Cref{fig:filteredgdp} illustrates the decomposition of US GDP into trend and cycle.

The parameter $\lambda$ penalizes variation in the trend: a larger value produces a smoother trend. Following \cite{Hodrick:1981aa}, we set $\lambda = 1,600$ for quarterly data.

\generatedfigure{filteredgdp.png}{US GDP Detrended with the Hodrick-Prescott Filter}{fig:filteredgdp}

Despite its widespread use, the HP filter has been criticized for its effects on measured business cycle properties. \cite{King:1993aa} find that it reduces standard deviations and substantially alters correlations between variables. They also argue that the conditions under which it minimizes the mean squared error are unlikely to hold in practice. Mechanical application can therefore remove variation normally associated with business cycles. Conversely, \cite{Cogley:1995aa} show that applying the filter to integrated time series can create apparent business cycle patterns that are absent from the underlying data. More recently, \cite{Phillips:2015aa} find that HP detrending generally fails to remove stochastic trends, contrary to a common assumption in applied macroeconomics.

The HP filter may also perform poorly during prolonged recessions.\footnote{A discussion of this issue in the economic blogosphere is summarized at \url{https://bruegel.org/2012/07/blogs-review-hp-filters-and-business-cycles/}.} During a deep slump, the estimated trend can flatten, reducing the measured deviation of output from trend. This can understate the depth or duration of a downturn and is a concern when interpreting the Great Recession. In our filtered US series, output remains below trend from 2008:4 to 2011:4. The National Bureau of Economic Research dates the end of the recession to June 2009, so the filtered series remains below trend beyond the officially dated end of the recession.